\begin{table}[!htbp]
\centering
\caption{School-ranking preferences by priority and continuation status: published student coordinates}
\label{tab:school-preference-groups-main}
\begin{tabular}{lcccc}
\toprule
 & \shortstack{School ends\\Non-prioritario} & \shortstack{School ends\\Prioritario} & \shortstack{Voluntary\\Non-prioritario} & \shortstack{Voluntary\\Prioritario} \\
\midrule
Math VA & 0.0253 (0.0029) & -0.0065 (0.0027) & 0.0457 (0.0068) & 0.0239 (0.0082) \\
Exam-taking VA & 0.0272 (0.0034) & -0.0441 (0.0025) & 0.0166 (0.0096) & -0.0020 (0.0093) \\
High-premium-field VA & 0.0615 (0.0024) & 0.0550 (0.0019) & 0.0560 (0.0064) & 0.0550 (0.0064) \\
High-premium-institution VA & -0.0196 (0.0048) & -0.0288 (0.0043) & 0.0452 (0.0117) & 0.0492 (0.0133) \\
log(1 + distance in km) & -0.1936 (0.0055) & -0.1716 (0.0040) & -0.2205 (0.0158) & -0.1645 (0.0165) \\
\midrule
Applicants & 87,407 & 155,028 & 15,363 & 13,799 \\
Student-school pairs & 338,722 & 525,303 & 57,911 & 46,638 \\
Destination schools & 2,484 & 2,497 & 2,261 & 2,304 \\
Log likelihood & -292,873.9970 & -396,102.8666 & -49,616.6263 & -35,336.8285 \\
\bottomrule
\end{tabular}
\par\smallskip
\begin{minipage}{0.98\linewidth}\footnotesize
\textit{Notes:} Separate joint rank-ordered logits of complete submitted school lists. Student-clustered standard errors are in parentheses. All four VA measures and distance enter jointly. VA retains the pooled model's equal-school normalization across 2,501 destination schools; it is not restandardized by group. School ends denotes an origin school without an SAE offering above grade 8; voluntary denotes an origin offering grade 9. These are application categories and do not require an eventual school move. Unresolved origins are excluded. Distance is straight-line kilometers and enters as $\log(1+d)$. The two distance specifications use identical samples within each group. Any missing regressor excludes the entire ranking; singleton lists provide no ordering information. Inference conditions on the existing EB VAs. Coefficients describe conditional ordering associations. Cross-group coefficient comparisons can also reflect differences in the unobserved utility-error scale.
\end{minipage}
\end{table}
